%%%%%%%%%%%%%%%%%%%%%%%%%%%%%%%%%%%%%%%%%
% Journal Article
% LaTeX Template
% Version 2.0 (February 7, 2023)
%
% This template originates from:
% https://www.LaTeXTemplates.com
%
% Author:
% Vel (vel@latextemplates.com)
%
% License:
% CC BY-NC-SA 4.0 (https://creativecommons.org/licenses/by-nc-sa/4.0/)
%
% NOTE: The bibliography needs to be compiled using the biber engine.
%
%%%%%%%%%%%%%%%%%%%%%%%%%%%%%%%%%%%%%%%%%

%----------------------------------------------------------------------------------------
%	PACKAGES AND OTHER DOCUMENT CONFIGURATIONS
%----------------------------------------------------------------------------------------

\documentclass[
	a4paper, % Paper size, use either a4paper or letterpaper
	10pt, % Default font size, can also use 11pt or 12pt, although this is not recommended
	unnumberedsections, % Comment to enable section numbering
	twoside, % Two side traditional mode where headers and footers change between odd and even pages, comment this option to make them fixed
]{LTJournalArticle}

\addbibresource{sample.bib} % BibLaTeX bibliography file
% Load float for placement specifier [H] which forces the float to be placed
% exactly where it appears in the source.
\usepackage{float}

\runninghead{Shortened Running Article Title} % A shortened article title to appear in the running head, leave this command empty for no running head

\footertext{\textit{Journal of Biological Sampling} (2024) 12:533-684} % Text to appear in the footer, leave this command empty for no footer text

\setcounter{page}{1} % The page number of the first page, set this to a higher number if the article is to be part of an issue or larger work

%----------------------------------------------------------------------------------------
%	TITLE SECTION
%----------------------------------------------------------------------------------------

\title{Practical Implementation and Evaluation of Ethereum's Proof-of-Stake Bouncing Attack and Patch on MOBS } % Article title, use manual lines breaks (\\) to beautify the layout

% Authors are listed in a comma-separated list with superscript numbers indicating affiliations
% \thanks{} is used for any text that should be placed in a footnote on the first page, such as the corresponding author's email, journal acceptance dates, a copyright/license notice, keywords, etc
\author{%
	Ricardo Loureiro\textsuperscript{1,3}, António Ravara\textsuperscript{1,3} and Simão Melo de Sousa\textsuperscript{2,3}
}

% Affiliations are output in the \date{} command
\date{\footnotesize\textsuperscript{\textbf{1}}FCT, NOVA\\ \textsuperscript{\textbf{2}} University of Algarve \\ \textsuperscript{\textbf{3}}NOVA-LINCS}

% Full-width abstract
\renewcommand{\maketitlehookd}{%
	\begin{abstract}
		\noindent With Ethereum going through \textit{The Merge} and moving from a Proof-of-Work to a Proof-of-Work consensus
							it has oppened itself to a new range of attacks. One of this attacks was a bouncing attack on liveness that
							prevented the blockchain from finalizing new blocks, the Ethereum community then responded with a patch to fix
							this vulnerability. An article analyzing this vulnerability and subsquent patch has found that there is a new
							vulnerability on the patched protocol that results in the protocol only having probablistic liveness. The article
							only provides a theorethical analysis and pseudo-code which motivated us to implement Ethereum as described, the attack
							and subsquent patch in MOBS. We therefore provide imperical evidence to complement the theoretical analysis supporting the
							claims made on the article.
	\end{abstract}
}

%----------------------------------------------------------------------------------------

\begin{document}

\maketitle % Output the title section

%----------------------------------------------------------------------------------------
%	ARTICLE CONTENTS
%----------------------------------------------------------------------------------------

\section{Introduction}

Consensus and blockchain protocols are usually described with pseudocode and model checked with idealized languages that 
not reflect the implementations. Since validating correctness is very costly, before planning an
implementation, developers should test their ideas and prototypes. There is a lack of methodologies to develop,
evaluate, tune and replicate these protocols, opening a need for extensible and modular tools for consensus analysis and
experimental labs for rapid prototyping.
With this in mind MOBS was developed with the aim to give the ability to test and validate these protocols under different conditions and settings by 
changing the execution parameters, aiming to catch vulnerabilities or even logic errors before deploying changes to these protocols.
Since verification of protocols is very costly, developers can use MOBS to get confidence in their implementations
from experimentations first.

Consensus protocols are inherently complex to design and implement correctly.
In blockchain systems, where financial transactions and dynamic behavior are the primary features, protocol correctness is essential,
as even minor errors can lead to substantial financial losses and system malfunctions. One illustration of this is the transition of
Ethereum from a Proof of Work to a Proof of Stake consensus. This transition exposed the protocol to vulnerabilities that could be
exploited by bouncing attacks on liveness.
The chain is unable to be finalized as a result of this type of attack, as the primary selected chain in the fork
choice rule is perpetually hopping between two alternative branches.

Liveness in a blockchain protocol is defined as the ability to always finalize new blocks. These blocks in the
Ethereum protocol are finalized through a process called Casper FFG (Friendly Finality Gadget), in combination with
the LMD GHOST (Latest Message Driven Greediest Heaviest Observed SubTree) fork choice rule.

The Bouncing Attack is an attack that prevents the chain from being finalized by making the main chain
selected in the fork choice rule continually bounce between two candidate chains. The attack leverages the fact that
candidate chains should start from the justified checkpoint with the highest epoch, and Byzantine nodes exploit this by
dividing honest validators opinions, justifying a new checkpoint after some honest validators already have cast their vote.

In~\cite{ethereum_analysis} this attack and with the Ethereum protocol are specified and analyzed along with the patch
implemented by the Ethereum community. In this article is also stated and formalized that the patch only provides probabilistic
liveness but no implementation is given. TODO Describe reason to do this implementation to support the article.

This document serves as a complement and provides more detailed information about our implementation and how it relates to the sdescription given in the article.

Here we show and detail our implementation in MOBS along with the re-creation of the Bouncing Attack, the patch proposed by the Ethereum
community and how this patch only provides probabilistic liveness.

%------------------------------------------------

\section{Implementation}

Our implementation was done in MOBS a simulator makes use of OCaml's \textit{modules} and \textit{functors} to
provide modularity and extensibility. MOBS adopts a \textit{Discrete-Event Simulation Model}
making the state of the system only change in discrete points in time when events occur.
These events are stored in a queue ordered with two main values:
\begin{itemize}
  \item \textbf{Timestamp:} This value dictates the order in which the events are
stored in the queue and is based on the simulator's internal clock. When getting
a new event the simulator will fetch from the queue the one with the smallest timestamp
and move its internal clock to match that event's timestamp.
  \item \textbf{Target:} The entity that should process this event.
\end{itemize}
Events are fetched from the queue until no more events remain or a predefined
stopping condition is reached.

For our implementation in MOBS we followed the description provided in pseudocode provided in~\cite{ethereum_analysis}. We focused on the
parts pertaining to the consensus mechanism and to allow a faithful replication of the results ignored some of the randomness mechanisms
that were described like choosing the member of the Proposers and the Validators committee. We also assumed that each attestation has a weight of 1
when the protocol was running for simplicity.


The first algorithm presented in the article~\ref{fig:article_algorithm1} represents the main execution loop. 
In our implementation~\ref{fig:mobs_receive_main} we continuously every node sends
a message to itself at the end of the main loop execution to execute the main code again. Also instead of extracting the logic of
\textit{prepareBlock()} and \textit{prepareAttestation()} to their own methods we keep them explicit in the main code.

\begin{figure}[H]
	\includegraphics[width=\linewidth]{article_algorithm1.png}
	\caption{Article algorithm 1}
	\label{fig:article_algorithm1}
\end{figure}

\begin{figure}[H]
	\includegraphics[width=\linewidth]{mobs_receive_main.png}
	\caption{MOBS receive_main function}
	\label{fig:mobs_receive_main}
\end{figure}

In our implementation of \textit{broadcastBlock()}~\ref{fig:article_algorithm2} we kept the same flow of the presented algorithm with the only
branches of execution being for the byzantine execution.

\begin{figure}[H]
	\includegraphics[width=\linewidth]{article_algorithm2.png}
	\caption{Article algorithm 2}
	\label{fig:article_algorithm2}
\end{figure}

For \textit{broadcastAttestation()}~\ref{fig:article_algorithm3} we ignored the Checkpoint vote logic since there was no description
on how a node should vote for the checkpoint and how it should use their stake on the vote. We did the assumption in our
execution that every node always voted with the same stake.

\begin{figure}[H]
	\includegraphics[width=\linewidth]{article_algorithm3.png}
	\caption{Article algorithm 3}
	\label{fig:article_algorithm3}
\end{figure}

For \textit{sync()}~\ref{fig:article_algorithm4} we did not explicitly implement this method. Since we determined the attestation and proposer
quorums to be static. The logic to call the \textit{justificationFinalization} is done in the \textit{receive\_main} method
between the \textit{prepareBlock()} and \textit{prepareAttestation()} methods.

\begin{figure}[H]
	\includegraphics[width=\linewidth]{article_algorithm4.png}
	\caption{Article algorithm 4}
	\label{fig:article_algorithm4}
\end{figure}

For \textit{syncBlock()}~\ref{fig:article_algorithm4} and \textit{syncAttestation()}~\ref{fig:article_algorithm6} we skipped over
the \textit{isValid()} method since there was no description on the implementation of this method, and we assumed that the byzantine
would not be interfering with these messages. 
Our implementation of \textit{syncBlock()}~\ref{fig:mobs_receive_propose} just updates the node's tree with the new block. To simplify
the process and avoid concurrency races we just update the whole tree with the received tree.
For our implementation of \textit{syncAttestation()} we follow the same logic and save the information of the attestation in a
hash table where each key is the epoch and the value is a list of attestation received. This will be used later for the
justification and finalization logic.
\begin{figure}[H]
	\includegraphics[width=\linewidth]{article_algorithm5.png}
	\caption{Article algorithm 5}
	\label{fig:article_algorithm5}
\end{figure}

\begin{figure}[H]
	\includegraphics[width=\linewidth]{article_algorithm6.png}
	\caption{Article algorithm 6}
	\label{fig:article_algorithm6}
\end{figure}

\begin{figure}[H]
	\includegraphics[width=\linewidth]{mobs_receive_propose.png}
	\caption{MOBS receive_propose}
	\label{fig:mobs_receive_propose}
\end{figure}

\begin{figure}[H]
	\includegraphics[width=\linewidth]{mobs_receive_attestation.png}
	\caption{MOBS receive_attestation}
	\label{fig:mobs_receive_attestation}
\end{figure}

For \textit{getHeadBlock()}~\ref{fig:article_algorithm7} we implemented the logic for this under the helper methods to navigate
and operate the n-ary tree that represents the blockchain. I would suggest consulting the repository and checking the \textit{deepest\_justified\_on\_path}
method since it would not be feasible to show the code through images because of the other helper methods that \textit{deepest\_justified\_on\_path} makes
use of.
\begin{figure}[H]
	\includegraphics[width=\linewidth]{article_algorithm7.png}
	\caption{Article algorithm 7}
	\label{fig:article_algorithm7}
\end{figure}

The \textit{getWeight()}~\ref{fig:article_algorithm8} was not implemented due to the fact that we assumed every node was voting with a stake of
1.
\begin{figure}[H]
	\includegraphics[width=\linewidth]{article_algorithm8.png}
	\caption{Article algorithm 8}
	\label{fig:article_algorithm8}
\end{figure}

Algorithm 9~\ref{fig:article_algorithm9} defines the consensus logic for the protocol. To implement this we followed the same logic as
presented in the article like shown in~\ref{fig:receive_justification_finalization}. When any node is finalized we propagate a
\textit{FinalizedNode} message, this acts as a logging mechanism and as away to converge every node to the same tree.
\begin{figure}[H]
	\includegraphics[width=\linewidth]{article_algorithm9.png}
	\caption{Article algorithm 9}
	\label{fig:article_algorithm9}
\end{figure}
\begin{figure}[H]
	\includegraphics[width=\linewidth]{mobs_receive_justification_finalization.png}
	\caption{MOBS receive_justification_finalization}
	\label{fig:mobs_receive_justification_finalization}
\end{figure}
%------------------------------------------------

\section{Results}

\begin{itemize}
	\item Basic Scenario
	\item Byzantine execution
	\item Byzantine Exceution with patch
	\item Exectuion showing the patch only provides probabilistic liveness
\end{itemize}

\section{Conclusion}
\begin{itemize}
	\item Relate implentation to the article (reforçar que é um trabalho empirico de confirmação das teses teoricas apresentadas no artigo, validação feita por simulação
	que implicou uma implementação dos algoritmos de pseudo codigo com decisoes estrategicas para capturar o essencial e deixar de lado o acessório)
	\item Future directions (implementação pode servir de guia para a implementação de outros protocolos de consensus semelhantes e de blockchains -> salientar também na introdução)
\end{itemize}

%----------------------------------------------------------------------------------------
%	 REFERENCES
%----------------------------------------------------------------------------------------

\printbibliography % Output the bibliography

%----------------------------------------------------------------------------------------

\end{document}
